\section{Uniform seventh, eighth and ninth polynomial minima}

\begin{finding}{R-11}{the seven-plane improvement in every matrix size}
Over every algebraically closed field of characteristic zero,
for every $n\geq4$ and for either polynomial family,
\[
 \mu_7(U_{n,2})=14,\qquad
 \mu_8(U_{n,2})=15,\qquad
 \mu_9(U_{n,2})=16.
\]
Consequently
\[
 \BR(\mathcal D_n),\ \BR(\mathcal P_n)
 \ \geq\ \binom n2^2+7\qquad(n\geq4).
\]
This is the consequence of the degree-two residual calculation;
it is not asserted to be the strongest known bound for general $n$.
\end{finding}
\begin{proof}
P-01 reduces the seventh minimum to seven distinct rectangular
quadratic minors or permanents. Let $G$ be the bipartite support
graph: its vertices are the appearing rows and columns and its
edges are the appearing matrix-entry variables. As in R-10,
the support flattening has a unit minor of order $|E(G)|$.
If $|E(G)|\geq14$, the desired lower bound follows.

Suppose therefore $|E(G)|\leq13$. Every edge belongs to one of
the seven selected rectangles. Every appearing row and column
has degree at least two. There are consequently at most six
row vertices and six column vertices. If there are at most four
of each, the tensor embeds, after relabelling and adjoining
unused coordinates, into the size-four situation. R-08 or R-10
then gives the lower bound fourteen. Such adjoining does not
change border rank: inclusion gives one inequality and a split
coordinate projection gives the reverse inequality.

In the remaining cases, transpose if necessary to put the larger
vertex count first. Write the row and column counts as $(r,c)$;
then $r\in\{5,6\}$ and $2\leq c\leq r$. A row is specified by
its nonempty neighbour subset of the $c$ columns. We enumerate
nondecreasing multisets of $r$ such subsets, each of size at
least two, with total size at most thirteen. We discard a
multiset if a column has degree below two, if fewer than seven
rectangles occur, or if some edge belongs to no rectangle.
Finally we identify column permutations, since the sorted row
multisets already identify row permutations. This covers every
possible support graph; it may retain graphs for which no
seven rectangles use every edge, which does not invalidate
the cover.

There are exactly thirty-four retained graph types. For each,
enumerate all seven-subsets of its rectangles whose union is
the whole graph, up to the graph's row/column automorphisms.
The exact cover has the following sizes:
\begin{center}
\begin{tabular}{c|r|r}
$(r,c)$&Support graphs&Seven-subset orbits\\\hline
$(5,2)$&1&4\\
$(5,3)$&6&853\\
$(5,4)$&19&116\\
$(5,5)$&1&6\\
$(6,2)$&1&20\\
$(6,3)$&5&27\\
$(6,4)$&1&1\\\hline
Total&34&1027
\end{tabular}
\end{center}
There are no retained types with $(r,c)=(6,5)$ or $(6,6)$.
The independent verifier enumerates all neighbour multisets
again, compares the complete canonical graph set, constructs
every vertex permutation preserving each graph, and checks
pairwise disjointness of the listed seven-subset orbits.
It then independently tests every seven-subset of that graph
to check exact coverage of those using every edge. A kernel
in the permutation action on rectangles is harmless: the
verifier uses the set of distinct induced actions.

For each family, first-factor Koszul minors give lower bound
fourteen for $1019$ of the $1027$ representatives. The remaining
eight have the following support-factor Koszul certificates.
Seven have thirteen variables; exterior degree one has pure
cost $12$, and the certified ranks are at least $158>13\cdot12$.
The final representative has twelve variables; exterior degree
three has pure cost $\binom{11}{3}=165$, and the ranks are at
least $2150$ for the determinant and $2147$ for the permanent,
both strictly greater than $13\cdot165=2145$.

The full graph cover, selected rectangles, orbit sizes and
minor indices are recorded in \path{uniform7_certificate.json}.
The independent \path{verify_uniform7.py} reconstructs every
selected Koszul entry from its tensor coefficient and exterior
insertion sign. The support-factor maps are decomposed into
distinct row/column torus weight blocks, whose ranks add.
Every listed minor is evaluated modulo the prime $1000003$,
with primality, coordinate ranges, distinct rows and columns,
block weights and the pure-tensor costs checked. Nonzero
residues certify nonzero integer minors and therefore the
required ranks in characteristic zero. The verifier imports
only the independently audited determinant and Koszul-entry
routines from \path{verify_det7_global.py}, not search code.
These finite certificates prove the lower bound in every
remaining case, hence $\mu_7\geq14$ for every $n\geq4$.

The anchored seven-edge construction of P-06 uses only four
rows and four columns and has border rank fourteen. It embeds
in every $n\geq4$, proving equality. R-07 then gives
$\mu_j\geq j+7$ for $j\geq7$, while the anchored eight- and
nine-edge constructions give equality for $j=8,9$.
The flattening image dimension is $\binom n2^2$, so R-01 gives
the asserted full-tensor bounds.
\end{proof}

\paragraph{Geometric dictionary.}
The finite support graph is an index for a coordinate subbundle
in the integral quadratic residual bundle. Split inclusions and
projections carry the small-support tensors between ambient
matrix sizes. P-01 supplies the torus specialization, and
exterior morphisms between locally free sheaves supply the
determinantal conditions. Graph automorphisms act through
permutation matrices, with determinant signs absorbed into
the first-factor basis. The published EGA, SGA~3 and SGA~6
locators are those of the common dictionary. Rank is certified
by the actual morphisms and their minors; no replacement by
a representation-ring class, purity weight or motive is used.
